\section{目标相对的恢复历史}\label{sec:recovery}

\subsection{目标相对状态与四类有向运动}

令 $x_t$ 表示相对于一个预先给定目标测量的一维状态，并规定恰在目标处 $x_t=0$。在吸烟应用中，可以取 $x_t=A_t-A_t^*$，其中 $A_t$ 是成瘾存量、$A_t^*$ 是当前治疗目标；也可以采用任何保留目标位置和状态排序的单调有界归一化。本文把 $x_t$ 作为原始对象；具体应用另行给出其运动定律。

当 $x>0$ 时，向 $0$ 移动表示从目标上方进行减量，而远离 $0$ 表示复发加剧；当 $x<0$ 时，远离 $0$ 表示进一步低于目标或出现失稳，而向 $0$ 移动则表示回到计划治疗水平。由此定义四个模块 $T_R$、$A_R$、$A_L$ 和 $T_L$，每个模块都可以具有自身的价值测度。

\subsection{共同记账与个体特定赋值}

一份恢复日记首先生成一个有向通行账户。账户记录各模块跨越的目标相对状态水平及其重复次数。随后，原始条件弱序在把非通行后果视为固定背景变量的条件下，对各模块的通行剖面赋予主观价值。对于同一个个体，在整个比较域内使用同一套内部一致的赋值体系；对于另一个个体，相同账户可以得到不同的赋值。

这里的条件化明确了通行次序比较时保持不变的对象。在给定运动定律下，绕行与直接减量可能带来不同的消费、健康或调整后果，因此不必要求二者在这些后果字面完全相同时还能同时可行。条件次序所隔离的是对通行账户的一个部分评价。在人为设计或自然控制的菜单中，如果被比较的一对方案具有相同的非通行后果，则观察到的选择会直接揭示这一通行次序；若这些后果不同，则具体应用需要把通行次序与其他后果模型结合起来。与程序效用和序列效用类似 \citep{LoewensteinPrelec1993,FreyBenzStutzer2004,FreyStutzer2005}，这一条件对象可以具有经济含义，而无需在每个应用中都能把通行与其他后果分别操纵。

支撑通行赋值属于个体内部对象；若要进行个体间比较，还需要额外的归一化。

\begin{figure}[pos=htbp,align=\centering]
 \centering
 \includegraphics[width=0.98\textwidth]{Figure_1_recovery_accounting.png}
 \caption{\textbf{目标相对恢复历史的有向记账。} 图~(a) 展示一条依次接近、跨过、离开并再次到达目标的路径。图~(b) 把非静止的单调片段分配到 $T_L$、$T_R$、$A_L$ 和 $A_R$ 四个模块。图~(c) 给出目标界面的记账约定。一次连续跨越在物理上是一段运动，但在账户中分为两个有向通行片段：进入目标的趋向目标 片段包含目标并记录一次目标达成，而离开目标的远离目标 片段不包含目标处的质量。}
 \label{fig:accounting}
\end{figure}

\subsection{目标界面}

只要一个单调片段跨越目标，通行账户就在 $0$ 处将其拆分。进入目标的趋向目标 片段以目标为终点并包含目标点；离开目标的远离目标 片段从目标出发但排除目标点。因此，一次连续跨越恰好记录一次目标达成。尽管如此，这两个有向片段由于属于不同模块，其穿孔部分仍可具有不同价值。

例如，从 $x=2$ 连续移动到 $x=-2$ 在物理上是一次减量过程。在账户中，它由两部分构成：先是从右侧的 $2$ 向 $0$ 的趋向目标 通行，其中包含一次右侧目标达成；随后是从 $0$ 向 $-2$ 的左侧 远离目标 通行，且不计入目标处质量。这一记账方式把“到达目标”的价值与“越过目标继续前进”的价值分离开来。

\subsection{三种历史实验}

固定右侧距离 $0<a<b<c$。考虑三条具有共同初始距离 $c$ 和终止距离 $b$ 的历史：
\[
 \gamma^D:c\longrightarrow b,
 \qquad
 \gamma^C:c\longrightarrow a\longrightarrow b,
 \qquad
 \gamma^A:c\longrightarrow0\longrightarrow b.
\]
直接历史 $\gamma^D$ 记录单调减量；绕行历史 $\gamma^C$ 先进一步改善至 $a$，随后复发回到 $b$；触及目标的历史 $\gamma^A$ 先到达目标，再返回 $b$。三条历史具有相同的有序端点，因此它们之间的差异能够隔离内部通行价值。

对于任意归一化、正则且相容的支撑赋值 $\Lambda$，下文的表示推出
\[
\begin{aligned}
 U_\Lambda(\gamma^C)-U_\Lambda(\gamma^D)
  &=2\kappa_{R,\Lambda}((a,b]),\\
 U_\Lambda(\gamma^A)-U_\Lambda(\gamma^D)
  &=2\kappa_{R,\Lambda}((0,b])+\alpha_{R,\Lambda}.
\end{aligned}
\]
在一个固定支撑赋值下，第一个差分隔离了目标之外的“减量--复发”绕行价值；第二个差分则隔离了额外的近目标通行对比以及目标达成价值。在整个支撑族上，这些基数分量仅被集合识别。第~\ref{sec:pha} 节和第~\ref{sec:implications} 节把同一几何转化为序数转移限制和有限数据锐界。仅依赖端点的子类预测 $\gamma^D\sim\gamma^C\sim\gamma^A$；三者之间任何严格排名都会否定“目标相对终止状态充分性”。

\subsection{匹配持续时间与流量暴露}\label{sec:matched}

即使一些熟悉的结果汇总指标已经被匹配，方案的通行记录仍可以不同。令一个\emph{几何方案}指定有序的有限步状态路径；具体应用还可以为该路径提供一个可行的计时遍历。下面的结果保持每个计时遍历本身不变，只把共同端点处的停留用作匹配工具。

\begin{proposition}[固定遍历下的匹配方案设计]\label{prop:matched-design}
固定 $0\le b<c$，并给定有限多个包含在 $[0,c]$ 中的有限步几何方案 $\pi^1,\ldots,\pi^K$，每个方案都从 $c$ 出发并终止于 $b$。对每个 $k$，令 $z_k:[0,\tau_k]\to[0,c]$（$\tau_k>0$）是方案 $\pi^k$ 在应用中可行的一个计时遍历。假设可行时间安排类对在每个 $z_k$ 两端任意加入停留封闭：对任意 $\delta,\beta\ge0$，先在 $c$ 停留 $\delta$，随后不重参数化地沿 $z_k$ 运行，最后在 $b$ 停留 $\beta$，仍然可行。令连续函数 $g:[0,c]\to\R$ 满足 $g(c)\ne g(b)$。则存在有限的 $\overline T$，使得对每个 $T>\overline T$，都存在由 $z_k$ 仅通过加入端点停留而得到的可行时间安排 $\widetilde\pi^k:[0,T]\to[0,c]$，并且
\[
 \int_0^T g\bigl(\widetilde\pi^k(t)\bigr)\,\dd t
\]
对所有 $k$ 都相同。每个方案的非静止遍历保持不变，因此其通行账户也保持不变。
\end{proposition}

\begin{proof}
令
\[
 m_k=\int_0^{\tau_k} g(z_k(t))\,\dd t,
 \qquad
 \Delta=g(c)-g(b),
 \qquad
 a_k=\frac{\tau_k g(b)-m_k}{\Delta}.
\]
选择
\[
 \overline T
 =2\max_{1\le k\le K}\{0,-a_k,\tau_k+a_k\}.
\]
对于任意 $T>\overline T$，定义初始与终止停留时间
\[
 \delta_k=\frac{T}{2}+a_k,
 \qquad
 \beta_k=\frac{T}{2}-\tau_k-a_k.
\]
根据 $\overline T$ 的定义，二者均严格为正，且 $\delta_k+\tau_k+\beta_k=T$。第 $k$ 个时间安排先在 $c$ 停留 $\delta_k$，然后不作重参数化地沿原计时遍历 $z_k$ 运行，最后在 $b$ 停留 $\beta_k$。其累计流量为
\[
 \delta_k g(c)+m_k+\beta_k g(b)
 =\frac{T}{2}\bigl(g(c)+g(b)\bigr),
\]
其中最后一个等式使用了 $a_k\Delta=\tau_k g(b)-m_k$。因此，所有时间安排具有相同的持续时间和累计流量。我们只增加了端点停留，所以根据假设可行性得到保留，同时通行账户不变。
\end{proof}

取 $g(x)=x$ 即可匹配累计目标相对暴露。特别地，上一小节中直接历史、目标外绕行历史以及触及目标历史的任意应用可行计时版本，都可以嵌入一个足够长的时间安排中，在不改变遍历速度的情况下实现共同端点、共同持续时间和共同累计暴露。条件 $g(c)\ne g(b)$ 是非退化条件：在两个端点之间重新分配停留时间会改变累计流量，却不改变持续时间和通行记录。作为特例，如果遍历时间本身可以自由缩放，则时间跨度限制消失：对任意给定的 $T>0$，可以先给有限个遍历安排足够短的共同遍历时间，然后再应用端点停留构造。该设计可以精确匹配一个选定的、能够区分端点的连续流量统计量；若还要匹配其他后果，则需要更多设计变量或一个明确的应用模型。

\subsection{比较域与可行选择}

表示域包含用来比较通行记录并识别其代数限制的有限步目标相对历史。每一个治疗问题都有自己的可行菜单，而它可能只是这一比较域中的一个很小子集。因此，循环是否可重复、是否存在反向机会以及是否允许触及目标，都是应用特定的可行性特征。

受控比较使通行次序具有经验解释。命题~\ref{prop:matched-design} 表明：一旦给定应用上可行的遍历，只要共同时间跨度足够长，就能仅通过加入端点停留，精确匹配持续时间和一个能够区分端点的累计流量后果，而无需改变遍历速度。其他后果可以通过更丰富的设计进行匹配、单独测量，或由提供可行菜单的具体应用模型处理。
